\section{The Smallest Thing That Runs}
\label{sec:ch02-the-smallest-thing-that-runs}

Start from an empty ASP.NET Core project rather than from the Hot Chocolate
template. The template writes files this book would have to print anyway, and
printing them is the point: every file this chapter needs is on the page,
complete, so you can type the service out of the book. Run the first command
inside a folder called \code{src}, because every service this book builds
lives beside this one under that name.

\begin{minted}{text}
dotnet new web -n Sessions
cd Sessions
dotnet add package HotChocolate.AspNetCore
dotnet add package HotChocolate.AspNetCore.CommandLine
dotnet add package HotChocolate.Types.Analyzers
\end{minted}

\index{HotChocolate.Types.Analyzers}%
Three packages, and only the first is obvious. \code{HotChocolate.AspNetCore}
is the server. The command line package adds subcommands that let the built
program export its own schema, which
section~\ref{sec:ch02-the-development-loop} turns into the loop you work in.
The analyzers package is a source generator, and
section~\ref{sec:ch02-where-the-registration-comes-from} is about what it
writes for you. That leaves the project file.

\begin{minted}[fontsize=\footnotesize]{xml}
<Project Sdk="Microsoft.NET.Sdk.Web">

  <PropertyGroup>
    <TargetFramework>net10.0</TargetFramework>
    <Nullable>enable</Nullable>
    <ImplicitUsings>enable</ImplicitUsings>
    <GenerateDocumentationFile>true</GenerateDocumentationFile>
    <NoWarn>$(NoWarn);CS1591</NoWarn>
  </PropertyGroup>

  <ItemGroup>
    <PackageReference Include="HotChocolate.AspNetCore" Version="16.6.1" />
    <PackageReference Include="HotChocolate.AspNetCore.CommandLine" Version="16.6.1" />
    <PackageReference Include="HotChocolate.Types.Analyzers" Version="16.6.1" />
  </ItemGroup>

</Project>
\end{minted}

\index{nullable reference types}%
Two of those properties are load-bearing and neither is here for style.
\code{Nullable} decides whether a field arrives in the schema as
\code{String} or as \code{String!}, which is the difference between a
client that copes with a missing value and one that does not. And
\code{GenerateDocumentationFile}, despite its name, decides whether your
XML comments reach the schema as descriptions or evaporate.
Section~\ref{sec:ch02-what-the-schema-says-about-your-types} shows both
happening. \code{NoWarn} then silences the warning about undocumented
public members that the second one would otherwise raise across the whole
project.

Replace the generated launch profile with this one, so that the port is yours
rather than whatever number the template picked. Every request in this chapter
goes to it.

\begin{minted}{json}
{
  "$schema": "https://json.schemastore.org/launchsettings.json",
  "profiles": {
    "http": {
      "commandName": "Project",
      "dotnetRunMessages": true,
      "launchBrowser": false,
      "applicationUrl": "http://localhost:5001",
      "environmentVariables": {
        "ASPNETCORE_ENVIRONMENT": "Development"
      }
    }
  }
}
\end{minted}

Now the program itself. Two files, and the whole service is in them. First
\code{Program.cs}, replacing what the template generated:

\begin{minted}{csharp}
var builder = WebApplication.CreateBuilder(args);

builder
    .AddGraphQL()
    .AddTypes();

var app = builder.Build();

app.MapGraphQL();

app.RunWithGraphQLCommands(args);
\end{minted}

Then \code{Query.cs}, which is the entire schema for now:

\begin{minted}{csharp}
namespace Sessions;

[QueryType]
public static class Query
{
    public static string Hello() => "world";
}
\end{minted}

\index{source generator!AddTypes}%
That is what the documentation tells you to write. It describes the template as
calling \enquote{a source-generated \code{AddTypes} method that registers all
types decorated with attributes like \code{[QueryType]} in the current
assembly}~\autocite{chillicreamgetstarted}. Build it and the compiler
disagrees.

\begin{minted}[fontsize=\footnotesize, breakanywhere]{text}
Program.cs(5,6): error CS0121: The call is ambiguous between the following methods or properties: 'Microsoft.Extensions.DependencyInjection.SchemaRequestExecutorBuilderExtensions.AddTypes(HotChocolate.Execution.Configuration.IRequestExecutorBuilder, params System.Type[])' and 'Microsoft.Extensions.DependencyInjection.SchemaRequestExecutorBuilderExtensions.AddTypes(HotChocolate.Execution.Configuration.IRequestExecutorBuilder, params HotChocolate.Types.ITypeDefinition[])'
\end{minted}

\code{AddTypes} is not the name of the generated method. It is the name of two
older overloads that each take a parameter array, and a call with no arguments
matches both of them equally well. What the generator wrote is called
\code{AddSessionsTypes}, after the assembly, which makes it a different
identifier in every project you will ever create. \code{Program.cs} again, with
the one line that changes marked:

\begin{minted}[highlightlines={5}]{csharp}
var builder = WebApplication.CreateBuilder(args);

builder
    .AddGraphQL()
    .AddSessionsTypes();

var app = builder.Build();

app.MapGraphQL();

app.RunWithGraphQLCommands(args);
\end{minted}

A compile error in a hello-world is a strange place to start a chapter, and it
is where the chapter starts because the error is the first thing this stack
tells you about itself. The method you call does not exist until you build, its
name came from the name of your assembly, and an attribute on a class is what
produced it. That is how the rest of the schema gets made too.

Run it and ask for the one field there is. Requests appear as raw HTTP
throughout this book, so no page depends on a tool you have not installed.

\begin{minted}{text}
POST /graphql HTTP/1.1
Host: localhost:5001
Content-Type: application/json

{"query":"{ hello }"}
\end{minted}

and back:

\begin{minted}{text}
HTTP/1.1 200 OK
Content-Type: application/graphql-response+json; charset=utf-8

{"data":{"hello":"world"}}
\end{minted}

\index{media type!application/graphql-response+json}%
Look at the media type in that response. It is not
\code{application/json}, and it is worth checking now what your own HTTP
client does with a type it has never met, because a client that unwraps
JSON by media type alone will hand you bytes.
Section~\ref{sec:ch02-the-development-loop} shows the same endpoint answering
with a status code you probably were not expecting either.
